% Generated draft: requires mathematical and source review.
\documentclass[11pt]{article}
\usepackage[margin=25mm]{geometry}
\usepackage{fontspec}
\setmainfont{Latin Modern Roman}
\usepackage{amsmath,amssymb,mathtools,unicode-math}
\setmathfont{Latin Modern Math}
\usepackage{enumitem,needspace,xurl,hyperref}
\setlength{\parindent}{0pt}
\setlength{\parskip}{5pt}
\setlength{\emergencystretch}{3em}
\newcommand{\N}{\mathbb{N}}
\newcommand{\Z}{\mathbb{Z}}
\newcommand{\Q}{\mathbb{Q}}
\newcommand{\R}{\mathbb{R}}
\newcommand{\C}{\mathbb{C}}
\newcommand{\F}{\mathbb{F}}
\newcommand{\A}{\mathbb{A}}
\newcommand{\PP}{\mathbb{P}}
\newcommand{\E}{\mathbb{E}}
\newcommand{\eps}{\varepsilon}
\newcommand{\dd}{\,\mathrm{d}}
\begin{document}
\section*{UnsolvedMath ID: 20001921}
\textbf{Source problem code:} catalogue entry
\subsection{Definitions and mathematical statement}
Let \(g_0\) be the canonical sectional-curvature-one metric on \(\R\mathrm P^n\), and put \(V_0=\operatorname{Vol}(\R\mathrm P^n,g_0)=\tfrac12\operatorname{Vol}(S^n)\). Its prime closed geodesics have length \(\pi\). The source asks whether every smooth Riemannian metric \(g\) with \(\operatorname{Vol}(g)=V_0\) has a nonconstant closed geodesic of length at most \(\pi\). Equivalently, for the shortest nonconstant closed-geodesic length \(L_{\min}(g)\), the scale-invariant question is
\[
L_{\min}(g)^n\leq\frac{\pi^n}{V_0}\operatorname{Vol}(g).
\]
The question does not require the geodesic to be noncontractible. For \(n\geq2\), the stronger noncontractible systole \(\operatorname{sys}_{\pi_1}(g)\) is the infimum of lengths of noncontractible loops and is realized by a closed geodesic. The case \(n=1\) is the circle-length normalization and is immediate; the source reports a positive answer for \(n=2\).

The supplied research aid gives the sharp conformal special case. If \(n\geq2\) and \(g=f^2g_0\), where \(f>0\) is smooth, then it reports
\[
\operatorname{sys}_{\pi_1}(f^2g_0)^n
\leq\frac{\pi^n}{V_0}\operatorname{Vol}(f^2g_0),
\]
with equality exactly when \(f\) is constant. Its stated method uses the invariant probability measure \(\mu\) on canonical projective lines, identified with two-planes in \(\R^{n+1}\), and the incidence average
\[
\int L_{f^2g_0}(\ell)\,\dd\mu(\ell)
=\frac\pi{V_0}\int_{\R\mathrm P^n}f\,\dd V_{g_0},
\]
followed by H\"older's inequality.

The aid also reports, for every \(n\geq2\), a smooth family \(h_t\) with pointwise volume form \(\dd V_{h_t}=\dd V_{g_0}\) and
\[
\int L_{h_t}(\ell)\,\dd\mu(\ell)\longrightarrow\infty.
\]
No formula for this family is supplied in the record. These \(\ell\) are the fixed canonical projective lines, not necessarily \(h_t\)-geodesics. Divergence of their average length does not bound the shortest \(h_t\)-geodesic from below and is not a counterexample to the source question. It only obstructs a naive extension of this particular average-length argument. The aid describes the unrestricted case \(n\geq3\) as still open.
\subsection{Short English statement}
At the canonical volume, must a metric on real projective space have a closed geodesic of length at most \(\pi\)? The supplied aid reports a sharp conformal systolic estimate. Its volume-preserving family only makes the average length of fixed canonical lines large and does not disprove the general claim.
\subsection{Sources}
\begin{itemize}
\item Question 4.11 of [Gromov2001], with the surface case cited as Proposition 4.10, in the AIM geodesics problem list: \url{https://aimath.org/pastworkshops/geodesicsproblems.pdf}. A full entry for [Gromov2001] is not supplied.
\item Local registry \texttt{data/unsolvedmath/problems.json}, permanent ID 20001921, code \texttt{catalogue entry}; canonical \texttt{aim-geometry-notes.json}, \texttt{notes[258]}, source item 4.11; updated 14 August 2026. The recorded verification reports that earlier PDF inspection of page 20 recovers a damaged label.
\item Supplied attempt 1: partial result, problem status at run open, rigor label proved, novelty confidence low; current status partially solved. The machine-generated aid is not peer reviewed and requires independent verification before citation; no detailed reference for the conformal estimate or family is supplied.
\end{itemize}
\end{document}
